\begin{table*}[!t]
\centering
\footnotesize
\renewcommand{\arraystretch}{0.82}
\setlength{\tabcolsep}{4.5pt}
\caption{电力市场结算诊断性审计：基于退化零申报反例 ($q \equiv 0$) 评估孤立不平衡现金流指标的有效性（基于 2017 年英国 Elexon SBP 全年出清时间线）。}
\label{tab:zero_bid_diagnostic}
\begin{threeparttable}
\begin{tabularx}{\textwidth}{p{2.5cm}p{3.5cm}>{\centering\arraybackslash}X>{\centering\arraybackslash}X>{\centering\arraybackslash}Xp{4.2cm}}
\toprule
\textbf{运行风电场} & \textbf{功率预测 / 出清申报} & \textbf{短缺清算罚款} & \textbf{单边平衡净现金流} & \textbf{实际决策成本} & \textbf{实证审计诊断结论} \\
\textbf{站点概况} & \textbf{策略类型} & \textbf{(\pounds/(台$\cdot$年))} & \textbf{(\pounds/(台$\cdot$年))} & \textbf{$\text{Cost}_{static}$ ($r=2$)} & \textbf{\& 方法学定性评价} \\
\midrule
\multirow{4}{=}{\textbf{Kelmarsh}\\(6台机组)} 
& \textbf{退化零申报 ($q \equiv 0$)} & \textbf{\pounds 0.00} & \textbf{+\pounds 223,882.15} & 0.30428 & \textbf{指标退化反例}：不进行预测即可凭空结算巨额账面现金流且罚款为零 \\
& 持续性预测基准 & \pounds 7,120.45 & +\pounds 5,890.12 & 0.09132 & 基础工程对照 \\
& 无约束 DLinear & \pounds 5,764.03 & +\pounds 6,584.74 & \textbf{0.08804} & 真实的决策损失最低；现金流数量级远低于退化策略 \\
& 校准经验分布先验 & \pounds 8,412.28 & +\pounds 3,485.65 & 0.09841 & 短缺罚款飙升；平衡收益显著萎缩 \\
\midrule
\multirow{4}{=}{\textbf{Hill of Towie}\\(21台机组)} 
& \textbf{退化零申报 ($q \equiv 0$)} & \textbf{\pounds 0.00} & \textbf{+\pounds 239,299.30} & 0.30473 & \textbf{指标退化反例}：仅凭全场入网电量即可结算出巨额账面现金流 \\
& 持续性预测基准 & \pounds 7,450.12 & +\pounds 9,120.34 & 0.08992 & 简单运行对照 \\
& 无约束 DLinear & \pounds 5,941.22 & +\pounds 8,678.09 & \textbf{0.08913} & 最小化真实的非负日前调度运行风险 \\
& 校准经验分布先验 & \pounds 5,702.68 & +\pounds 12,174.95 & 0.09503 & 现金流虚假激增 \pounds 3.6k，但实际决策损失恶化 $+6.5\%$ \\
\midrule
\multirow{4}{=}{\textbf{Penmanshiel}\\(14台机组)} 
& \textbf{退化零申报 ($q \equiv 0$)} & \textbf{\pounds 0.00} & \textbf{+\pounds 243,449.80} & 0.33689 & \textbf{指标退化反例}：证实单边现金流与预测精度本质脱钩 \\
& 持续性预测基准 & \pounds 7,080.20 & +\pounds 6,150.30 & \textbf{0.09586} & 复杂地形下的强劲线性基准 \\
& 无约束 DLinear & \pounds 6,658.86 & +\pounds 6,377.76 & 0.09852 & 小样本基准模型 \\
& 校准经验曲线 ($\lambda=1.0$) & \pounds 6,682.66 & +\pounds 6,382.43 & \textbf{0.09726} & 真实物理先验成功降低实际调度决策风险 \\
\bottomrule
\end{tabularx}
\begin{tablenotes}[flushleft]
\footnotesize
\item \textbf{注记}：(1) 不平衡现金流作为简化 BSC 价格驱动的诊断模型计算为 $\sum (y_t - q_t) \times P_{\text{imb}, t}$，其中 $q_t = \hat{y}_t$ 为简化核算模型中的前瞻合同或申报头寸。本诊断模型抽象掉了双边合同量通知 (ECVN)、平衡机制动作与残余现金流再分配。对于退化零申报策略 ($q \equiv 0$)，$y_t - q_t = y_t \ge 0$，故短缺恒为零，所有未申报上网电量均按 $P_{\text{imb}}$ 计入溢出收入。(2) 脱离远期合同持仓（$P_{\text{F}} \cdot q_t$）而孤立评估未对冲平衡现金流，不能作为衡量预测商业价值的独立指标。(3) 非对称决策成本 ($r=2, q=1/3$) 准确惩罚了零申报策略的高额低估风险（$\text{Cost}_{\text{static}} = \bar{y} \approx 0.30$--$0.34$，对应各站测试集平均容量系数）。(4) DLinear 与经验先验行报告 5 个相同随机种子 (42, 123, 456, 789, 1024) 上的均值，与表~\ref{tab:controlled_ablation} 保持完全一致。
\end{tablenotes}
\end{threeparttable}
\end{table*}
